% Folder 156-9, batch 05 — archivists' pages 81–100.
% Pass: Opus 5.5 (claude-opus-5-5), 2026-10-03 — first pass, unchecked against the pages by a human.
%
% Notation fixed for this batch: $\mathrel{|\circ|}$ for his relation written
% « |o| » between elements, parts or subsets (his « disposition »);
% $\Sigma_E$ the system attached to a disposition $(E,|\circ|_E)$, $\sigma_E$
% the canonical map $E \to \Sigma_E$; $\mathfrak{P}(E)$ for the set of parts;
% $\complement$ for the complement he writes as a C; $\operatorname{Sup}^{\circ}$,
% $\operatorname{Cosup}^{\circ}$ as he writes them, with the small circle.
\documentclass[11pt,a4paper]{article}
\input{../preamble/grothendieck}

\folder{156-9}
\batch{5}
\pages{81}{100}
\foldertitle{[Chapitre] IX et IX bis. [Ateliers] : notes manuscrites (05-15/07/1986).}
\dating{1986}
\watermark{Édition de démonstration}

\begin{document}

\page{81}
\note{pagination de l'auteur : p. 78. L'argument vient de la batch précédente : la condition a), la relation $\mathrel{|\circ|}$, $\Sigma$, $\Sigma_E$ et $\sigma$ y sont déjà introduits.}

b) $\forall\, S \in \Sigma$, posant $\beta(S) = \{x \in E \mid \sigma(x) \leq S\}$, on a
\struck{$S = \operatorname{Sup}(\{\sigma x \mid x \in$}
\[
  S = \operatorname{Sup}_{x \in \beta(S)} \sigma x
\]
(Donnons aussi la formulation \struck{symétrique} \add{b'), p.~83}).
\note{« p.~83 » est la pagination de l'auteur, soit la p.~86 des archivistes.}

\marginal{Cette condition b) \uncertain{signifie aussi} que « $\sigma(E)$ est \struck{partie} \add{« dense »} de $\Sigma$ », i.e. $\forall\, S$, $S = \operatorname{Sup}_{\sigma(x) \leq S} \sigma(x)$ \uncertain{\ill{}}.}

De plus, quand ces conditions sont satisfaites, \uncertain{la} \ill{} $\alpha$ est uniquement déterminée. \textbf{NB} Explicitons la formule pour $\alpha$, et pour l'application inverse $\beta$.

\emph{Dém.} La \uncertain{nécessité} de a) b) évidente, prouvons la suffisance. Utilisant a), on \uncertain{va} construire une application
\[
  \alpha : \Sigma_E \longrightarrow \Sigma,
\]
en posant pour tout $A \in \Sigma_E$
\[
  \alpha(A) = \operatorname{Sup}_{x \in A} \sigma(x).
\]
Il est clair que cette application est croissante, et par b) elle est surjective, \struck{\ill{}} en \uncertain{vérifiant} que $\beta(S)$ dans b) est $\in \Sigma_E$. Soit $S' = \complement S$, et considérons $\beta(S')$, il suffit de voir que
\[
  \beta(S) = \operatorname{Cosup}^{\circ}(\beta(S'))
\]
i.e. que pour $x \in E$, on a
\[
  x \in \beta(S) \iff x \mathrel{|\circ|} \beta(S') \quad \text{i.e.} \quad \sigma(x) \mathrel{|\circ|} \underbrace{\sigma(\beta(S'))}_{S'_0}
\]
i.e. $\sigma(x) \leq S$ (par a)).

Mais $\sigma(x) \mathrel{|\circ|} S'_0 \iff \sigma(x) \mathrel{|\circ|} \operatorname{Sup} S'_0$, et comme $\operatorname{Sup} S'_0 = S'$, cela équivaut à $\sigma(x) \mathrel{|\circ|} S'$ i.e. $\sigma(x) \leq S$, ok.

\page{82}
\note{p.~79 de l'auteur.}

En fait, on \add{\uncertain{dans}} construit une application
\[
  \Sigma \xrightarrow{\ \beta\ } \Sigma_E
\]
\uncertain{réciproquement} croissante, telles que
\[
  \alpha\beta = \mathrm{id}_{\Sigma} .
\]
\note{le $\alpha$ de $\alpha\beta$ est récrit sur une lettre noircie.}

De plus, on vérifie \uncertain{de suite} qu'elle est compatible : C. Pour démontrer ce dernier, il suffit de voir que
\[
  \beta\alpha = \mathrm{id}_{\Sigma_E} ,
\]
\note{après $\beta\alpha$, quelques lettres noircies et illisibles.}
car alors $\alpha$ et $\beta$ sont des \struck{applications} \add{iso. d'un ens. ordonné sur l'autre}, \struck{\uncertain{i.e.} l'une \ill{}} \add{inverses l'une de l'autre}, et (comme on \uncertain{voit} sur $\beta$) compatibles : C.

Soit donc $A \in \Sigma_E$, $A' =$ \struck{\ill{}} $\complement A$. Alors les $\sigma(x)$ et $\sigma(x')$ ($x \in A$, $x' \in A'$) \add{\uncertain{par ex.}} sont disjoints, \add{donc} \uncertain{leurs Sup} \struck{\ill{}} $\alpha(A)$, $\alpha(A')$ le sont \uncertain{aussi}. \struck{Considérons}
\struck{$\beta\alpha(A) = \{x \in E \mid \sigma(x) \leq \alpha(A)\}$} \struck{Considérons} il s'ensuit \uncertain{d'où}
\[
  \alpha(A) \leq \complement\alpha(A') ,
\]
i.e.
\[
  \beta(\alpha(A)) \leq \beta\complement\alpha(A') = \complement\beta(\alpha(A')) , \quad \text{i.e.} \ldots
\]
Mais cela implique
\[
  \beta\alpha(A) \mathrel{|\circ|} \beta\alpha(A') ,
  \qquad
  \beta\alpha(A) \geq A , \quad \beta\alpha(A') \geq A' ,
\]
\note{les deux inégalités sont écrites verticalement, sous chacun des deux termes.}
$A = \beta\alpha(A)$, $A' = \beta\alpha(A')$, par le

\emph{Lemme.} Soit $\Sigma$ un système ($\Sigma_E$ ou l'\uncertain{analogue}), $A \leq B$, $A' \leq B'$ avec $B \mathrel{|\circ|} B'$, et $A' = \complement A$. Alors $A = B$, $A' = B'$.

En effet, \uncertain{comme on a toujours} \ill{}
\[
  A \leq B \leq \complement B' \leq \complement A' = A ,
\]
donc \uncertain{toutes ces $\leq$ sont des égalités}.

\page{83}
\note{p.~80 de l'auteur.}

Il reste : voir l'unicité de $\alpha$, qui est immédiate, vu que \struck{tout} $A \in \Sigma_E$ est le sup des $\sigma_E(x)$ qu'il majore (i.e. telles que $x \in A$ \ldots), \struck{et $\alpha$ doit \uncertain{être}} d'où forcément
\[
  \alpha(A) = \operatorname{Sup}_{x \in A} \underbrace{\alpha(\sigma_E(x))}_{\sigma(x)} = \operatorname{Sup}_{x \in A} \sigma(x) ,
\]
\uncertain{cqfd}.

\struck{\emph{Corollaire.} Soit $(E, \mathrel{|\circ|})$ un \uncertain{dispositif} \ill{}, $E_0$ une partie très \ill{} de $E$}
\note{ces deux lignes sont barrées de traits obliques.}

\emph{Corollaire !} Soient $(E, \mathrel{|\circ|}_E)$ une disposition, $E'$ une partie de $E$, $\mathrel{|\circ|}_{E'}$ la relation induite. Conditions équivalentes :
\begin{enumerate}
  \item[(i)] L'image $\sigma_E(E')$ de $E'$ dans $\Sigma_E$ est dense, i.e. tout $S \in \Sigma_E$ est le sup des éléments de $\sigma_E(E')$ qu'il majore.
  \item[(ii)] \struck{\ill{}} L'application
  \[
    S \longmapsto S \cap E' , \qquad \Sigma_E \longrightarrow \mathfrak{P}(E')
  \]
  \struck{applique $\Sigma_E$ dans $\Sigma_{E'}$, et \uncertain{induit} \ill{}} \add{est} \struck{surjection} (\textbf{NB} elle induit \uncertain{donc} une iso $\varphi : \Sigma_E \to \Sigma_{E'}$). \struck{(NB. \ill{} une iso)}
  \item[(iii)] L'application
  \[
    \psi : \Sigma_{E'} \longrightarrow \Sigma_E , \qquad S' \longmapsto \operatorname{Sup}^{\circ}_E(S')
  \]
  est \uncertain{surjection} (NB \uncertain{en ce sens} $=$ une iso)
  \item[(iv)] $E'$ est \emph{dense} dans $E$ pour la relation \add{de pré}ordre canonique de $E$ (induite de celle de $\Sigma_E$ via $\sigma_E$).
\end{enumerate}
\marginal{\uncertain{permutés} \struck{\ill{}} avec (v) $\longleftarrow$}
\note{la marge renvoie par une flèche à la condition (iv).}

\page{84}
\note{p.~81 de l'auteur.}

Sous ces conditions, les applications $\varphi$, $\psi$ sont inverses l'une de l'autre, et \ill{} des \uncertain{opérations}.

\add{(v)} $\forall\, x, y \in E$, si $\forall\, x', y'$ avec $x' \leq x$, $y' \leq y$, \uncertain{on a} $x' \mathrel{|\circ|} y'$, alors $x \mathrel{|\circ|} y$. (\uncertain{La} relation $\leq$ est la relation d'ordre canonique)
\note{la condition (v), d'abord numérotée autrement (numéro noirci), est rattachée par une accolade et une flèche au haut de la page.}

\emph{Dém.} (i) $\iff$ (iv), \add{évident}, compte tenu que (iv) signifie aussi que $\sigma(E')$ est dense dans $\sigma(E)$, pour l'ordre induit par $\Sigma_E$. (Attention, « dense » ici \uncertain{ne veut pas dire} \ill{} de la relation d'ordre, qui \struck{est pas} \uncertain{n'a} rien à voir avec la top. logique associée \uncertain{à} l'ordre ; mais \uncertain{on a} une \uncertain{propriété} de transitivité de la densité, grâce à l'associativité des sup.)

(i) $\Longrightarrow$ (v) est connu une \uncertain{fois} que $\sigma(E') \subset \sigma(E) \subset \Sigma$, et immédiat, car dans $\Sigma$, si $x = \operatorname{Sup} x_i$, $y = \operatorname{Sup} y_j$, $x_i \mathrel{|\circ|} y_j$ $\forall\, i, j$, alors $x \mathrel{|\circ|} y$.

(v) $\Longrightarrow$ (iv) \struck{Soit $S \in \Sigma_E$,} Dans \uncertain{la situation} \ill{} $\Sigma$ avec $\sigma(E') \subset \sigma(E) \subset \Sigma$. Soit $\xi \in \sigma(E)$. \struck{Disons que} Considérons \struck{$S = \sigma(x)$ ($\in \sigma(E)$)} l'ens.\ $S'$ des $\xi' \in \sigma(E')$ tels que $\xi' \leq \xi$. Je dis que
\[
  \xi = \operatorname{Sup} S' .
\]
\struck{\ill{}} Mais \uncertain{(cf. lemme)} cela signifie aussi que pour tout $\eta \in \sigma(E)$, on a $S' \mathrel{|\circ|} \eta \Rightarrow \xi \mathrel{|\circ|} \eta$. Considérons donc $J' = \{\eta' \in \sigma(E') \mid \eta' \leq \eta\}$. On a donc $\xi' \mathrel{|\circ|} \eta'$ pour $\xi' \in S'$, $\eta' \in J'$. Par \uncertain{hyp.}, cela implique $\xi \mathrel{|\circ|} \eta$, ok.

\page{85}
\note{p.~82 de l'auteur.}

\emph{Lemme.} Soient $\Sigma$ un système, $E$ une partie \uncertain{dense} (cf.~\ldots), $S'$ une partie de $\Sigma$, \uncertain{formée} d'éléments majorés par un $\xi \in \Sigma$, \struck{donc} i.e. $\operatorname{Sup} S' \leq \xi$. Pour que $\xi = \operatorname{Sup} S'$, il faut et suffit que pour tout $y \in E$, \struck{tel que}
\[
  \text{\struck{$S' \mathrel{|\circ|} y$}} \quad S' \mathrel{|\circ|} y \Longrightarrow \xi \mathrel{|\circ|} y .
\]

\emph{Dém du Lemme.} On sait déjà que si $\operatorname{Sup} S' = \xi$, alors $\forall\, y \in \Sigma$, on a
\[
  S' \mathrel{|\circ|} y \iff \xi \mathrel{|\circ|} y .
\]
En sens inverse, \struck{supposons que $\operatorname{Sup}(S') = \xi$ i.e. $\xi \leq \operatorname{Sup} S'$,} \struck{$\complement\xi \geq \ldots$}
\[
  \forall\, y \in E, \quad
  \underbrace{S' \mathrel{|\circ|} y}_{\text{i.e. } y \leq \operatorname{Inf}_{\xi' \in S'} \complement(\xi')}
  \Longrightarrow
  \underbrace{\xi \mathrel{|\circ|} y}_{\text{i.e. } y \leq \complement(\xi)}
\]
À cause de la densité, on en conclut
\[
  \operatorname{Inf}_{\xi' \in S'} \complement(\xi') \leq \complement(\xi) ,
\]
d'où
\[
  \xi \leq \operatorname{Sup}_{\xi' \in S'} \xi' ,
\]
\uncertain{cqfd}.

Ainsi, les conditions (i) (iv) (v) de la prop.\ sont équivalentes. Compte tenu de prop.~1, la condition (i) exprime que l'application canonique $E' \subset E \to \Sigma_E$ permet d'identifier $\Sigma_E$ à $\Sigma_{E'}$. Cette identification est donnée par (\uncertain{ii})

\page{86}
\note{p.~83 de l'auteur ; la phrase reprend la fin de la page précédente.}

(cf.\ construction de l'application $\beta$ dans prop.~2), et en sens inverse par (iii) (application $\alpha$ de prop.~2). Donc, (i) \uncertain{(ou les équivalentes} (iv), (v)) implique (ii), (iii) \uncertain{sauf pour} \ill{} (i.e. \ill{} de \uncertain{opérations}). Il reste : prouver que (ii) ou (iii) implique \add{(déjà)} (i).

Mais que (iii) (et \uncertain{simplement} la \emph{surjectivité} de $\psi$) implique (i), est clair. Donc (iii) équivaut aux \uncertain{autres} conditions. Reste : voir que (ii) implique (i). Cela équivaut au corollaire suivant, appliqué à $E' \to \Sigma_E$ :

\emph{Corollaire 2.} Soient $E$ une disposition, $\Sigma$ un système, $\sigma : E \to \Sigma$ une application\ldots\ Alors la condition b) de la prop.~2 est équivalente (moyennant a)) à la condition suivante :

b') L'application $S \mapsto \beta(S)$ de $\Sigma$ dans $\mathfrak{P}(E)$ applique $\Sigma$ dans $\Sigma_E$, et est \struck{bijective} \add{injective}.

\uncertain{On a vu} dans la \uncertain{démonstration} que b) $\Rightarrow$ b'). Prouvons que b') $\Rightarrow$ b). Considérons, pour $S \in \Sigma$ fixé,
\[
  S' = \operatorname{Sup} \sigma(\beta(S)) \leq S .
\]
Il est clair que $\beta(S) = \beta(S')$, donc $S = S'$, puisque $\beta$ est injectif, cqfd.

\page{87}
\note{p.~84 de l'auteur. Les corollaires 3 et 4 n'apparaissent pas : la numérotation saute du corollaire 2 au corollaire 5.}

\emph{Corollaire 5.} Soient $(E, \leq, \mathrel{|\circ|})$ une disposition \uncertain{prédonnée}, et soit $E'$ une partie de $E$ qui est « dense » dans le sens suivant :
\[
  \forall\, x, y \in E, \quad E'_{\leq x} \mathrel{|\circ|} E'_{\leq y} \Longrightarrow x \mathrel{|\circ|} y .
\]
\note{avant $E'_{\leq x}$, un signe noirci et illisible.}
Alors
\[
  \Sigma_{E'} \underset{\varphi}{\overset{\psi}{\rightleftarrows}} \Sigma_E
\]
\note{les deux flèches portent le signe $\sim$.}
iso.\ réciproques l'un de l'autre, donnés par
\[
  \psi(A') = \operatorname{Sup}^{\circ}_E(A') , \qquad \varphi(A) = A \cap E' .
\]

On applique cor.~1, critère (v), en \uncertain{notant} que la condition de densité ici est plus forte a priori que celle dans (v).

\section{\uncertain{8 bis.} Correspondances entre \add{dispositions et entre} systèmes}
\note{titre souligné dans le manuscrit ; le numéro qui le précède est d'une lecture douteuse.}

\add{\uncertain{Dic.}} \emph{Analogie.} Soient $E$, $E'$ deux ensembles, \struck{\ill{}} \uncertain{vus comme} dispositions triviales. Une « corr » (ou relation \add{$R$}) entre $E$, $E'$ peut \uncertain{s'}interpréter comme une application
\[
  \alpha : E \longrightarrow \mathfrak{P}(E') \qquad (x \longmapsto R(x))
\]

\page{88}
\note{p.~85 de l'auteur.}

et elle se prolonge \uncertain{aussitôt} en une application
\[
  \alpha_{*} \text{ ou } \overline{\alpha} : \mathfrak{P}(E) \longrightarrow \mathfrak{P}(E') , \qquad A \longmapsto \bigcup_{x \in A} \alpha(x) .
\]
Cette application $\overline{\alpha}$ commute aux Sup (\uncertain{donc} est croissante), et inversement, \struck{\ill{}} toute application $\varphi : \mathfrak{P}(E) \to \mathfrak{P}(E')$ qui a cette propriété, \uncertain{provient} \uncertain{d'une} unique $\alpha$, sous la forme $\alpha_{*}$, puisque $\forall\, A \in \mathfrak{P}(E)$, on a
\[
  A = \operatorname{Sup}_{x \in A} \{x\} \quad \text{d'où} \quad \varphi(A) = \operatorname{Sup}_{x \in A} \underbrace{\varphi(\{x\})}_{\text{\uncertain{déf.} } \alpha(x)} .
\]

\marginal{J'écris $\overline{\alpha}$ au lieu de $\alpha_{*}$, \uncertain{pour} \uncertain{la} \uncertain{suite} : \ill{} dis \uncertain{aussi} une application \ill{} $\mathfrak{P}(E)$, \ill{} comme \uncertain{flèche} \ill{} \ill{}}
\note{note marginale écrite en oblique dans l'angle supérieur gauche, en grande partie illisible.}

Nous allons interpréter l'application correspondante \uncertain{symétrique}, ou transposée, $\beta$ de $\alpha$, \uncertain{ou} \uncertain{plutôt}
\[
  \beta_{*} \text{ ou } \overline{\beta} : \mathfrak{P}(E') \longrightarrow \mathfrak{P}(E) ,
\]
en \struck{\ill{}} termes des opérations $\mathfrak{P}(E)$, $\mathfrak{P}(E')$, en notant que
\[
  \forall\, A \in \mathfrak{P}(E),\ A' \in \mathfrak{P}(E') \quad \text{on a} \quad
  \alpha_{*}(A) \mathrel{|\circ|} A' \iff A \mathrel{|\circ|} \beta_{*}(A') .
\]
\note{dans le membre de droite, le $\mathrel{|\circ|}$ est écrit sans le petit rond : $A \mathbin{\|} \beta_{*}(A')$.}

En effet, le premier \uncertain{membre} s'écrit
\[
  \text{\struck{$\forall\, x \in$}}\ \forall\, x' \in A', \ \text{on a } x' \notin \overline{\alpha}(A) \ \text{ i.e. } \nexists\, x \in A
\]
avec $x' \in \alpha(x)$ i.e. $(x, x') \in R$
\[
  \text{on a encore} \quad \forall\, x \in A,\ x' \in A' \quad \text{on a } (x, x') \notin R \ \text{ i.e. } (x, x') \in \overline{R}
\]
\note{la phrase se poursuit sur la page suivante.}

\page{89}
\note{p.~86 de l'auteur.}

et \uncertain{le} \uncertain{second} \uncertain{traduit} le \ill{} \uncertain{interprétation} pour le 2e \uncertain{terme} \uncertain{des} \uncertain{deux}.

Soient maintenant $(E, \mathrel{|\circ|}_E)$, $(E', \mathrel{|\circ|}_{E'})$ deux dispositions, et considérons deux correspondances
\[
  \alpha : E \to \mathfrak{P}(E') \ \text{ou}\ \overline{\alpha} : \mathfrak{P}(E) \to \mathfrak{P}(E') ,
  \qquad
  \beta : E' \to \mathfrak{P}(E) \ \text{ou}\ \overline{\beta} : \mathfrak{P}(E') \to \mathfrak{P}(E) .
\]

\emph{Prop.~3} \struck{\emph{Lemme}} Conditions équivalentes :
\begin{enumerate}
  \item[(i)] $\forall\, x \in E$, $x' \in E'$, on a
  \[
    x \mathrel{|\circ|} \beta(x') \iff \alpha(x) \mathrel{|\circ|} x' .
  \]
  \item[(ii)] $\forall\, A \subset E$, $A' \subset E'$, on a
  \[
    A \mathrel{|\circ|} \overline{\beta}(A') \iff \overline{\alpha}(A) \mathrel{|\circ|} A' .
  \]
  \item[(iii)] Posant, pour $A \in \Sigma_E$, $A' \in \Sigma_{E'}$,
  \[
    \alpha_{*}(A) = \operatorname{Sup}^{\circ}_{E'}(\overline{\alpha}(A)) , \qquad
    \beta_{*}(A') = \operatorname{Sup}^{\circ}_{E}(\overline{\beta}(A')) ,
  \]
  \note{après « $\in$ », avant $\Sigma_E$, un mot noirci et illisible ; les lettres $A$, $A'$ sont récrites sur d'autres.}
  on a aussi équivalence
  \[
    \text{\struck{$\alpha_{*}(A)$}}\ A \mathrel{|\circ|} \beta_{*}(A') \iff \alpha_{*}(A) \mathrel{|\circ|} A' .
  \]
\end{enumerate}
\struck{De plus, \ill{}} \struck{plus bas relations, pour $A \in \mathfrak{P}(E)$, $A' \in \mathfrak{P}(E')$ :}
\[
  \text{\struck{$\alpha_{*}(\operatorname{Sup}^{\circ}_E(A)) = \operatorname{Sup}^{\circ}_{E'}(\overline{\alpha}(A))$,}}
  \quad
  \text{\struck{$\beta_{*}(\operatorname{Sup}^{\circ}_E(A')) = \operatorname{Sup}^{\circ}_E(\overline{\beta}(B))$}}
\]
\note{ce passage est barré, et une flèche venue de la marge (« transposées ») le renvoie plus bas.}

\textbf{Cor.} Ces conditions \uncertain{impliquent} que \uncertain{les} applications $\alpha_{*}$, $\beta_{*}$
\[
  \Sigma_E \underset{\beta_{*}}{\overset{\alpha_{*}}{\rightleftarrows}} \Sigma_{E'}
\]
\uncertain{définissent} \ill{} \uncertain{pour} $\alpha$ et $\beta$, et en particulier un \uncertain{choix} \uncertain{global} pour \uncertain{\ill{}} \uncertain{remplaçons} $\alpha$ et $\beta$ par $\alpha'$, $\beta'$ définis par
\struck{$\alpha(x) = \operatorname{Sup}^{\circ}_{E'} \alpha(x)$}
\note{le bas de la page est barré de deux longs traits obliques.}

\marginal{Définition. On dit dans ce cas \uncertain{que} $\alpha$ \uncertain{et} $\beta$ sont \uncertain{transposées} \ill{} \uncertain{l'une de l'autre}. \ldots\ \uncertain{dès} \uncertain{que}}
\note{note oblique dans la marge gauche, d'une lecture très incertaine.}

\page{90}
\note{p.~87 de l'auteur. Le haut de la page, jusqu'à « transposées $\alpha_1$, $\alpha_2$ de $\beta$ », est barré d'un long trait oblique, qui continue celui de la page précédente.}

\[
  \alpha'(x) = \operatorname{Sup}^{\circ}_{E'} \alpha(x) , \qquad
  \beta'(x) = \operatorname{Sup}^{\circ}_{E} \beta(x) .
\]
De plus, pour \struck{\ill{}} \add{donnés} $\alpha$ (ou $\alpha_{*}$ donné), \struck{et plus,} \add{ou $\alpha'$,} \struck{il existe une seule $\beta_{*}$ qui} si $\beta_1$, $\beta_2$ sont deux transposées de $\alpha$ \ldots\ ($\beta'_1 = \beta'_2$, $\beta_{1*} = \beta_{2*}$), et \uncertain{ainsi} \add{des côtés} \struck{\ill{}} $\alpha$ si \ill{} deux transposées $\alpha_1$, $\alpha_2$ de $\beta$.

\emph{Dém.} Comme \uncertain{c'est} évident \uncertain{que} (ii) $\Rightarrow$ (i), en faisant $A = \{x\}$, $A' = \{x'\}$, et (i) $\Rightarrow$ (ii), on note que
\[
  A \mathrel{|\circ|} \overline{\beta}(A') \iff x \mathrel{|\circ|} \overline{\beta}(x') \quad \forall\, x \in A,\ x' \in A'
\]
\[
  \alpha(A) \mathrel{|\circ|} A' \iff \alpha(x) \mathrel{|\circ|} x' .
\]
\struck{De plus, $A$, $A' \in \Sigma_E$, lorsque \uncertain{tous}}
\note{« De plus » est surchargé d'un ajout au-dessus de la ligne, « en plus » ; la ligne entière est barrée.}

\struck{Et donc (ii) \uncertain{passe}} Notons que les deux \uncertain{membres} de (ii) ne changent pas quand on \uncertain{remplace} $\overline{\beta}$ par le \uncertain{Sup} des termes dans la relation $\mathrel{|\circ|}$, \ldots\ les deux, par \uncertain{leur} $\operatorname{Sup}^{\circ}$\uncertain{respectif}. Quand $A \in \Sigma_E$, $A' \in \Sigma_{E'}$, cela montre que \ldots\ cette équivalence \struck{signifie}, elle \uncertain{dans} (iii) donc équivalente. \add{a) On va prouver b) plus bas.}

(ii) $\Longrightarrow$ (iii). \struck{En particulier, prouvons que} (iii) $\Rightarrow$ (ii) (\uncertain{p.~\ldots}). On écrit la relation (ii) \uncertain{précédemment} sous la forme équivalente :
\note{la suite de cette démonstration est sur la page suivante.}

\page{91}
\note{p.~88 de l'auteur.}

\[
  (\ast) \qquad
  \underbrace{\operatorname{Sup}^{\circ}_E(A)}_{\overline{A}} \mathrel{|\circ|} \operatorname{Sup}^{\circ}_E(\overline{\beta}(A'))
  \iff
  \operatorname{Sup}^{\circ}_{E'}(\overline{\alpha}(A)) \mathrel{|\circ|} \underbrace{\operatorname{Sup}^{\circ} A'}_{\overline{A'}}
\]
\struck{Si on prouve que \ill{}}
\[
  \text{\struck{$\operatorname{Sup}^{\circ}_E(\overline{\beta}(A')) = \operatorname{Sup}^{\circ}_E$}}
\]
(On \uncertain{démontre} (iii) b) \uncertain{comme} \uncertain{sup}\ldots) \struck{\uncertain{Si} on prouve}
\[
  (\ast\ast) \quad
  \left\{
  \begin{aligned}
    \operatorname{Sup}^{\circ}_E(\overline{\beta}(A')) &= \operatorname{Sup}^{\circ}_E(\overline{\beta}(\overline{A'})) \quad \bigl(\overset{\text{déf}}{=} \beta_{*}(\overline{A'})\bigr) \\
    \operatorname{Sup}^{\circ}_{E'}(\overline{\alpha}(A)) &= \operatorname{Sup}^{\circ}_{E'}(\overline{\alpha}(\overline{A})) \quad \bigl(\overset{\text{déf}}{=} \alpha_{*}(\overline{A})\bigr)
  \end{aligned}
  \right.
\]
on en conclut que $(\ast)$ s'écrit aussi
\[
  \overline{A} \mathrel{|\circ|} \beta_{*}(\overline{A'}) \iff \alpha_{*}(\overline{A}) \mathrel{|\circ|} \overline{A'}
\]
ce qui \uncertain{n'est} autre que (iii). Donc il \struck{\uncertain{faut}} \add{suffit} prouver (\add{\uncertain{puisque} (ii) \uncertain{implique},} les deux relations \uncertain{($\ast\ast$)} --- et il suffit \struck{de \uncertain{prouver}}, \uncertain{par} raison de symétrie) de prouver la seconde :

\emph{Corollaire !} Sous les conditions du lemme 1 (\uncertain{ou} plutôt sous la forme (ii)), si $A \in \mathfrak{P}(E)$, $A' \in \mathfrak{P}(E')$, $\overline{A} = \operatorname{Sup}^{\circ}_E(A)$, $\overline{A'} = \operatorname{Sup}^{\circ}_{E'}(A')$, \ldots\ les relations \uncertain{ci-dessus}.

\emph{Dém.} Comme $A \subset \overline{A}$, \ldots\ $\overline{\alpha}(A) \subset \overline{\alpha}(\overline{A})$, d'où
\[
  \operatorname{Sup}^{\circ}_{E'}(\overline{\alpha}(A)) \subset \operatorname{Sup}^{\circ}_{E'}(\overline{\alpha}(\overline{A})) ,
\]
il \uncertain{reste} à prouver \uncertain{l'inclusion} \uncertain{inverse}. Je \uncertain{montre} \uncertain{qu'on} a $=$ \ldots\ que $\forall\, A' \in \Sigma_{E'}$, \ldots
\[
  A' \mathrel{|\circ|} \operatorname{Sup}^{\circ}_{E'}(\overline{\alpha}(A)) \iff A' \mathrel{|\circ|} \operatorname{Sup}^{\circ}_{E'}(\overline{\alpha}(\overline{A}))
\]
\note{les deux $\operatorname{Sup}^{\circ}_{E'}$ de cette ligne sont récrits sur des lettres noircies, et le second membre porte en dessous une première version raturée, illisible.}

\page{92}
\note{p.~89 de l'auteur. Le haut de la page, jusqu'à la flèche marquée « ? », est barré de deux traits obliques.}

\struck{Or par (iii), le dernier membre de l'implication équivaut à : $\beta_{*}(A') \mathrel{|\circ|} A$, alors que le premier équivaut à : $A' \mathrel{|\circ|} \overline{\alpha}(A)$. Reste donc à montrer que}
\[
  \text{\struck{$A' \mathrel{|\circ|} \overline{\alpha}(A) \overset{?}{\Longrightarrow} \beta_{*}(A') \mathrel{|\circ|} A$}}
\]
Or par (ii), les deux membres équivalent resp.\ à $\beta_{*}(A') \mathrel{|\circ|} A$, et $\beta_{*}(A') \mathrel{|\circ|} \struck{\ill{}}\ \overline{A}$, \uncertain{où} $\overline{A} = \operatorname{Sup} A$, lesquelles sont \uncertain{bien} équivalentes, qed.

En particulier, \struck{\uncertain{faisant}} \uncertain{avec} $A = \{x\}$, \ldots\ \uncertain{on a}
\[
  \operatorname{Sup}^{\circ}_{E'} \alpha(x) = \operatorname{Sup}^{\circ}_{E'} \overline{\alpha}\bigl(\operatorname{Sup}^{\circ}_E(x)\bigr) \overset{\text{déf}}{=} \alpha_{*}\bigl(\operatorname{Sup}^{\circ}_E(x)\bigr) ,
\]
i.e.
\[
  \boxed{\ \alpha_{*}\bigl(\operatorname{Sup}^{\circ}_E(x)\bigr) = \underbrace{\operatorname{Sup}^{\circ}_{E'}(\alpha(x))}_{\overset{\text{déf}}{=}\ \alpha'(x)}\ }
\]
\note{le $\operatorname{Sup}^{\circ}_E$ du membre de gauche est récrit sur un premier essai.}

donc \struck{la composition de $\alpha$ avec $\operatorname{Sup}^{\circ}_{E'}$ \ldots}

$\alpha_{*}$ et $\alpha'$ se déterminent mutuellement par cette formule et par la formule, \uncertain{évidente} quelle que soit $\alpha : E \to \mathfrak{P}(E')$ (\uncertain{quitte} à définir $\alpha'$ par composition \uncertain{de $\alpha$} avec $\operatorname{Sup}^{\circ}_{E'}$) :
\[
  \boxed{\ \alpha_{*}(S) = \operatorname{Sup}_{x \in S}^{(\Sigma_{E'})} \alpha'(x) \quad (= \alpha'_{*}(S))\ }
\]

Les conditions de \uncertain{Prop.~3} \uncertain{sur} $(\alpha, \beta)$ ne dépendant que de $\alpha'$, $\beta'$ \uncertain{associés}, \struck{desquelles \uncertain{définissent}} \add{\uncertain{les}} \uncertain{donnant} $\alpha'_{*} = \alpha_{*}$, $\beta'_{*} = \beta_{*}$,

\page{93}
\note{p.~90 de l'auteur.}

\[
  \left\{
  \begin{aligned}
    \alpha' &: E \to \Sigma_{E'} \\
    \beta' &: E' \to \Sigma_E
  \end{aligned}
  \right.
  \qquad
  \left\{
  \begin{aligned}
    \alpha_{*} &: \Sigma_E \to \Sigma_{E'} \\
    \beta_{*} &: \Sigma_{E'} \to \Sigma_E
  \end{aligned}
  \right.
\]
\marginal{NB $\alpha \longmapsto (\alpha' \leftrightarrow \alpha_{*})$, \quad $\beta \longmapsto (\beta' \leftrightarrow \beta_{*})$}

Ceci vu, je dis que la paire $\alpha$ (ou $\alpha'$, $\alpha_{*}$) \uncertain{étant} donnée, $\beta_{*}$ (ou $\beta'$) est déterminé de façon unique, par la condition de la proposition. En effet, $\beta_{*}(S')$ est connu quand on connaît les $S \in \Sigma_E$ tels que $S \mathrel{|\circ|} \beta_{*}(S')$, \uncertain{lesquels} \uncertain{sont} \uncertain{ceux} pour lesquels on a $\alpha_{*}(S) \mathrel{|\circ|} S'$ !

Je \uncertain{montre} : voir quelles sont les applications
\[
  \alpha : E \longrightarrow \Sigma_{E'} ,
\]
--- ou ce qui revient au même, les correspondances
\[
  \alpha_{*} : \Sigma_E \to \Sigma_{E'} , \qquad \alpha_{*}(S) = \operatorname{Sup}^{\Sigma_{E'}}_{x \in S} \alpha(x)
\]
pour lesquelles il existe une $\beta$
\[
  \beta : E' \longrightarrow \Sigma_E ,
\]
\struck{\ill{} ce qui revient} \uncertain{\ill{}}
\[
  \beta_{*} : \Sigma_{E'} \to \Sigma_E , \qquad \beta_{*}(S') = \operatorname{Sup}^{\Sigma_E}_{x' \in S'} \beta(x')
\]
transposée.

On \uncertain{a} \uncertain{quand} \ldots\ $=$ partie des $\alpha : E \to \mathfrak{P}(E')$ \uncertain{ou} $E \to \Sigma_{E'}$.

\marginal{\emph{Déf.} On dit que $\alpha$, $\beta$ \uncertain{sont} \uncertain{transposées} \uncertain{l'une} \uncertain{de} \uncertain{l'autre} \ldots\ $\alpha$, $\beta$ \ldots\ \ill{}}
\note{note oblique dans la marge gauche, en grande partie illisible ; elle reprend la définition esquissée en marge de la page 89.}

\page{94}
\note{p.~91 de l'auteur.}

\emph{Proposition 4.} Soit $\alpha : E \to \mathfrak{P}(E')$ une correspondance, \struck{\uncertain{considérons}} posons
\[
  \alpha' : E \to \Sigma_{E'} , \quad \alpha'(x) \overset{\text{déf}}{=} \operatorname{Sup}^{\circ}_{E'}(\alpha(x))
\]
\[
  \alpha_{*} : \Sigma_E \to \Sigma_{E'} , \quad
  \alpha_{*}(S) = \operatorname{Sup}^{\circ}_{E'}(\overline{\alpha}(S)) = \operatorname{Sup}^{\Sigma_{E'}}_{x \in S} \alpha'(x) .
\]
\note{au-dessus de cette dernière formule, une première version de la définition de $\alpha_{*}$ est raturée de traits serrés et illisible.}

Conditions équivalentes :
\begin{enumerate}
  \item[(i)] $\alpha$ admet une transposée
  \[
    \beta : E' \longrightarrow \mathfrak{P}(E)
  \]
  --- ou, ce qui revient au même, une transposée $\beta' : E' \to \Sigma_E$, ou $\beta_{*} : \Sigma_{E'} \to \Sigma_E$.
  \item[(ii)] a) $\alpha_{*}$ commute aux Sup quelconques ;
  b) $\forall\, x \in E$, $\alpha_{*}\bigl(\operatorname{Sup}^{\circ}_E(x)\bigr)$
  \[
    \text{(i.e. } \operatorname{Sup}^{\circ}_{E'}\bigl(\alpha(\operatorname{Sup}^{\circ}_E(x))\bigr)\text{)} = \operatorname{Sup}^{\circ}_{E'}(\alpha(x))
  \]
  \item[(iii)] On a la formule, pour \struck{$A$} $A \in \mathfrak{P}(E)$,
  \[
    \alpha_{*}\bigl(\underbrace{\operatorname{Sup}^{\circ}_E(A)}_{\overline{A}}\bigr) = \operatorname{Sup}^{\circ}_{E'}(\overline{\alpha}(A))
  \]
  (i.e.\ $\overline{\alpha}(\overline{A})$ et $\overline{\alpha}(A)$ ont \uncertain{même} Sup dans $E'$)
  \item[(iv)] $\forall\, S' \in \Sigma_{E'}$, l'ensemble
  \[
    S = \{x \in E \mid \alpha(x) \subset S'\}
  \]
  dans $E$ est dans $\Sigma_E$.
  \item[(v)] $\forall\, x' \in E'$, l'ensemble
  \[
    \beta(x') = \{x \in E \mid \alpha(x) \mathrel{|\circ|} x'\}
  \]
  est dans $\Sigma_E$.
\end{enumerate}
\note{les numéros (ii), (iii), (iv) ont été corrigés sur la page ; une accolade avec une flèche lie (ii) a) et b), et une longue flèche descend de (ii) jusqu'au bas de la page.}

\page{95}
\note{p.~92 de l'auteur.}

\emph{Dém.} Prouvons d'abord
\[
  \text{(i)} \Longrightarrow \text{(ii)} \Longrightarrow \text{(iii)} \Longrightarrow \text{(iv)} \Longrightarrow \text{(i)} ,
\]
ce qui établira l'équivalence de toutes les conditions sauf (v).
\note{le premier (ii) de la chaîne est récrit sur un chiffre noirci ; le (v) qui suit « sauf » l'est aussi.}

(i) $\Longrightarrow$ (ii) a été vu.

(ii) $\Longrightarrow$ (iii). \struck{Soit $T' = \complement S'$. On a alors, par} Soit $\overline{S} = \operatorname{Sup}^{\circ}_E S \supset S$. On a \struck{$x \in E$}, par (ii), $\overline{\alpha}(S)$ et $\overline{\alpha}(\overline{S})$ \struck{$S \subset \overline{S}$, \uncertain{déf} et par $S'$, $x \in S \iff \alpha(x) \subset S'$} ont \uncertain{même} \uncertain{Sup} dans $E'$, \uncertain{donc} \uncertain{contenu} \uncertain{dans} $S'$
\[
  \overline{\alpha}(S) \subset S' \Longrightarrow \overline{\alpha}(\overline{S}) \subset S' ,
\]
\note{ce passage passe de la numérotation (ii) $\Rightarrow$ (iii) à l'argument de (iii) $\Rightarrow$ (iv) sans transition nette ; transcrit dans l'ordre de la page.}
et \uncertain{comme} $\overline{\alpha}$ est un \uncertain{morphisme}, $\overline{\alpha}$ \ldots\ donc $\overline{S} \subset S$ par définition de $S$, donc $S = \overline{S}$, i.e.\ $S \in \Sigma_E$.

(iii) $\Longrightarrow$ (iv) \uncertain{de même}, prenant $S' = \operatorname{Cosup}^{\circ}\{x'\}$.

(iv) $\Longrightarrow$ (i), on \uncertain{trouve} la transposée
\[
  \beta(x') = \complement\, \gamma(x') \quad \ldots
\]
En effet, \struck{la relation} \add{\uncertain{des}} qui \uncertain{\ill{}}, pour $x \in E$, $x' \in E'$
\[
  x \mathrel{|\circ|} \beta(x') \iff \alpha(x) \mathrel{|\circ|} x'
\]
\[
  \text{i.e.} \quad \text{\struck{$\operatorname{Sup}^{\circ}_E(x) \subset$}} \quad x \in \complement\,\beta(x') = \gamma(x')
\]
ce qui \uncertain{n'est} autre que la définition de $\gamma(x')$.

Je dis que les conditions \add{équiv.\ (i) : (iv)} \uncertain{impliquent} \add{b) est clair par (ii) et \uncertain{prouvons} a)} (v). En effet, \struck{\ill{}}, soit $(S_i)_{i \in I}$ \uncertain{des} $S_i \in \Sigma_E$,
\[
  S = \operatorname{Sup}^{\Sigma_E} S_i = \operatorname{Sup}^{\circ}_E\Bigl(\bigcup S_i\Bigr) . \quad \text{On a} \ldots
\]
donc
\[
  \alpha_{*}(S) = \alpha_{*}\Bigl(\operatorname{Sup}^{\circ}_E\Bigl(\bigcup S_i\Bigr)\Bigr)
  = \operatorname{Sup}^{\circ}_{E'}\Bigl(\underbrace{\overline{\alpha}\Bigl(\bigcup S_i\Bigr)}_{\bigcup_i \overline{\alpha}(S_i)}\Bigr)
\]
\note{la lettre $\gamma$, introduite ici sans définition visible, désigne manifestement l'ensemble $\{x \in E \mid \alpha(x) \subset \operatorname{Cosup}^{\circ}\{x'\}\}$ de (iv) ; l'ajout « b) est clair par (ii) et prouvons a) » semble viser (ii) plutôt que (v) : la page ne le précise pas.}

\page{96}
\note{p.~93 de l'auteur.}

\[
  = \operatorname{Sup}^{\circ}_{E'} \bigcup_i \underbrace{\operatorname{Sup}^{\circ}_{E'}(\overline{\alpha}(S_i))}_{\alpha_{*}(S_i)} = \operatorname{Sup}_i \alpha_{*}(S_i) , \quad \text{ok}
\]

Je reste : prouver que (v) implique les conditions (i)\ldots\ \uncertain{par} \uncertain{exemple} (ii) a). \uncertain{Prouvons} la formule
\[
  \alpha_{*}\bigl(\operatorname{Sup}^{\circ}_E(A)\bigr) \overset{?}{=} \operatorname{Sup}^{\circ}_{E'}(\overline{\alpha}(A)) .
\]
Or on a
\[
  \operatorname{Sup}^{\circ}_E(A) = \operatorname{Sup}^{\Sigma_E}_{x \in A} \sigma_E(x)
\]
d'où, par (v),
\[
  \alpha_{*}\bigl(\operatorname{Sup}^{\circ}_E(A)\bigr) = \operatorname{Sup}^{\Sigma_{E'}}_{x \in A} \alpha_{*}(\sigma_E(x))
\]
et d'autre part
\[
  \operatorname{Sup}^{\circ}_{E'}\bigl(\underbrace{\overline{\alpha}(A)}_{\bigcup_{x \in A} \alpha(x)}\bigr) = \operatorname{Sup}^{\Sigma_{E'}}_{x \in A} \operatorname{Sup}^{\circ}_{E'} \alpha(x)
\]
et on \uncertain{a} \uncertain{dit} que, pour \uncertain{tout} \uncertain{$x$} \uncertain{seul},
\[
  \alpha_{*}\bigl(\operatorname{Sup}^{\circ}_E(x)\bigr) = \operatorname{Sup}^{\circ}_{E'}(\alpha(x))
\]
c'est (v~b)), que j'ai dû \uncertain{ajouter} \uncertain{après} \uncertain{coup}.
\note{la condition (v) de la page 94 n'a pas de partie b) visible : celle-ci est, semble-t-il, l'ajout que la phrase mentionne.}

\struck{La condition (v~b) \uncertain{exprime} \uncertain{aussi} que}

\emph{Cor.} \add{\uncertain{Scholie}} \struck{Soient} $\Sigma = \Sigma_{\underline{E}}$, $\Sigma' = \Sigma_{\underline{E}'}$, soit
\[
  \operatorname{Corr}(\Sigma, \Sigma')
\]
l'ensemble des applications
\[
  \alpha_{*} : \Sigma \longrightarrow \Sigma'
\]
qui commutent aux Sup (quelconques), \add{de $\Sigma$ : $\Sigma'$} (\uncertain{appelons} les « correspondances » \struck{\ill{} et} \uncertain{entre} \uncertain{elles} sont des \uncertain{connaissances}), et
\[
  \operatorname{Corr}(\underline{E}, \underline{E}')
\]
\note{ici $\underline{E}$, $\underline{E}'$ (soulignés par l'auteur) désignent les dispositions. La phrase se poursuit sur la page suivante.}

\page{97}
\note{p.~94 de l'auteur.}

l'ensemble des \struck{applications} correspondances entre $\underline{E}$, $\underline{E}'$
\[
  \alpha : E \longrightarrow \text{\struck{$\Sigma_{E'}$}}\ \mathfrak{P}(E')
\]
\note{le $E$ source est récrit sur un $\underline{E}$ ; avant $\mathfrak{P}(E')$, un $\Sigma_{E'}$ raturé de traits serrés.}
ayant les propriétés
\begin{enumerate}
  \item[a)] $\forall\, x \in E$, $\alpha(x) \in \Sigma_{E'}$
  \item[b)] $\alpha$ satisfait \struck{aux \uncertain{\ill{}}} \add{\uncertain{aux}} conditions équivalentes de la prop.~4, dont la plus explicite est \struck{\uncertain{\ill{}}} (iv), \struck{\ill{}}
\end{enumerate}

Soit enfin
\[
  \operatorname{Corr}^{*}(\underline{E}, \underline{E}')
\]
l'ensemble des \struck{applications parties} \add{corr.\ ensemblistes $R$} entre \struck{$\alpha^{*} : E \to \mathfrak{P}(E')$} $E$ et $E'$ ($R \subset E \times E'$) telles que
\[
  \left\{
  \begin{aligned}
    &\forall\, x \in E, \ R(x) \in \Sigma_{E'} \\
    &\forall\, x' \in E', \ {}^{t}R(x') \in \Sigma_E .
  \end{aligned}
  \right.
\]

On a donc des \struck{correspondances} \add{applications} bijectives \uncertain{canoniques}
\[
  \operatorname{Corr}(\Sigma, \Sigma') \xrightarrow[\sim]{\ \varphi\ } \operatorname{Corr}(\underline{E}, \underline{E}') \xrightarrow[\sim]{\ \psi\ } \operatorname{Corr}^{*}(\underline{E}, \underline{E}')
\]
\note{le second $\operatorname{Corr}^{*}$ est souligné dans le manuscrit.}
Obtenues ainsi :
\[
  \left\{
  \begin{aligned}
    \varphi(\alpha_{*}) &= \alpha_{*} \circ \sigma_E : \ E \xrightarrow{\ \sigma_E\ } \Sigma_E \xrightarrow{\ \alpha_{*}\ } \Sigma_{E'} \\
    \psi(\alpha) &= \{(x, x') \in E \times E' \mid \alpha(x) \mathrel{|\circ|} x'\} \\
    &\phantom{=}\ \text{i.e. } x' \in \operatorname{Cosup}^{\circ}_{E'}(\alpha(x))
  \end{aligned}
  \right.
\]
\note{le dernier argument, en bout de ligne, est coupé au bord de la feuille.}

(Ainsi $\psi(\alpha)$ n'est pas le graphe de la correspondance $\alpha : E \to$ \struck{$\mathfrak{P}(E')$} $\Sigma_{E'} \subset \mathfrak{P}(E')$, mais de la correspondance déduite en composant avec $\complement_{\Sigma'}$ :
\[
  E \xrightarrow{\ \alpha\ } \Sigma_{E'} \xrightarrow{\ \complement_{E'}\ } \Sigma_{E'} \subset \mathfrak{P}(E') \text{).}
\]

\page{98}
\note{p.~95 de l'auteur.}

Je reste seulement à prouver que si $\alpha_{*}$ est dans $\operatorname{Corr}(\Sigma, \Sigma')$, alors $\alpha = \varphi(\alpha_{*})$ satisfait aux conditions de prop.~4 et redonne $\alpha_{*}$, par
\[
  \alpha_{*}(S) = \operatorname{Sup}^{\Sigma_{E'}}_{x \in S} \underbrace{\alpha(x)}_{\overset{\text{déf}}{=}\ \alpha_{*}(\operatorname{Sup}^{\circ}_E(x))} \qquad (S \in \Sigma_E) .
\]
Cette formule \struck{\ill{}} provient de la formule
\[
  S = \operatorname{Sup}_{x \in S} \operatorname{Sup}^{\circ}_E(x)
\]
et de la commutation de $\alpha_{*}$ avec Sup.

Ainsi, il reste à prouver (iv~b) pour $\alpha$ :
\[
  \overline{\alpha}\bigl(\text{\struck{$\operatorname{Sup}^{\circ}_E$}}\ \sigma_E(x)\bigr) \text{ \struck{$\in$} et } \alpha(x) \text{ ont même } \operatorname{Sup}^{\circ}_{E'} .
\]
Or si $y \in \text{\struck{\ill{}}}\ \sigma_E(x)$, on a $\sigma_E(y) \leq \sigma_E(x)$
\struck{($\operatorname{Sup}^{\circ}_E(y) \leq \operatorname{Sup}^{\circ}_E(x)$)},
d'où \struck{$\alpha_{*}(\ldots)$}
\[
  \alpha_{*}\sigma_E(y) \leq \alpha_{*}\sigma_E(x)
  \quad \text{i.e.} \quad
  \underbrace{\alpha(y) \leq \alpha(x)}_{\text{dans } \Sigma_{E'}} ,
\]
\note{sous $\alpha_{*}\sigma_E(y)$, une accolade porte « $\alpha(y)$ », barré.}
donc \struck{les $\overline{\alpha}$}
\[
  \operatorname{Sup}^{\circ}_{E'}\bigl(\overline{\alpha}(\sigma_E(x))\bigr) = \ldots\ \alpha(x)
\]
\struck{$\alpha_{*}(\sigma_E(x))$, $\overline{\alpha}(\sigma_E(x))$ et $\alpha(x)$} qed.
\note{le second membre de la dernière égalité est raturé de traits serrés, illisible sauf le $\alpha(x)$ final ; deux lignes en dessous sont barrées.}

\uncertain{11.7.} On va faire une \emph{catégorie} avec les dispositions, en prenant comme morphismes
\[
  \underline{E} \longrightarrow \underline{E}'
\]
les correspondances qui satisfont a) b) de la p.~94, et la composition (si $\underline{E} \xrightarrow{\ \alpha\ } \underline{E}' \xrightarrow{\ \beta\ } \underline{E}''$)
\[
  \beta \underset{\text{dis}}{\circ} \alpha = \text{\struck{$\beta \underset{\text{ens}}{\circ} \alpha$}}\ \operatorname{Sup}^{\circ}_{E''}\bigl(\beta \underset{\text{ens}}{\circ} \alpha\bigr) .
\]
\note{« p.~94 » est la pagination de l'auteur : c'est la page 97 ci-dessus. Le numéro « 11.7 » en marge est d'une lecture douteuse.}
Cette

\page{99}
\note{p.~96 de l'auteur.}

composition est associative, avec \struck{les} \add{pour} identités
\[
  \text{\struck{\ill{}}}\ \sigma_{\underline{E}} = \operatorname{Sup}^{\circ}_{\underline{E}} \circ \mathrm{id}_{\underline{E}} .
\]
La catégorie des dispositions et correspondances \add{entre} \struck{de} telles, est équivalente à la catégorie des dispositions \uncertain{complètes} (ou « \uncertain{systèmes} » ?), où on prend comme morphismes les \emph{applications} croissantes qui commutent aux Sup quelconques (\uncertain{appelées} \uncertain{aussi} « correspondances » entre \uncertain{les} \uncertain{systèmes}). Et cette catégorie est équivalente à la catégorie opposée, par une « autodualité » (\uncertain{involution}) qui est l'identité sur les objets, et sur les morphismes est le passage à la transposée.

[Analogie avec les espaces de Hilbert, \uncertain{plus généralement} espaces vectoriels munis d'une autodualité\ldots, et les opérateurs adjoints\ldots]
\note{le mot « systèmes » de la première parenthèse est souligné d'un trait ondulé ; les crochets sont de l'auteur.}

\emph{Questions ouvertes}, que je \uncertain{n'ai} pas \uncertain{tirées} au clair, \uncertain{signalées} par le cas des dispositions triviales.
\note{« Questions ouvertes » est souligné d'un double trait ondulé, et une flèche d'une autre encre y pointe depuis la marge. La page s'arrête là : les questions annoncées ne figurent pas dans cette batch.}

\page{100}
\note{p.~97 de l'auteur.}

\emph{Dans le cas} où les correspondances en dispositions sont les corr.\ ensemblistes, elles forment l'ensemble $\mathfrak{P}(E \times E')$, qui est lui-même un \struck{disposition} système, associé à une disposition triviale $E \times E'$. Il serait tentant donc, dans le cas général, de \uncertain{considérer} les correspondances \uncertain{entre} $\underline{E}$, $\underline{E}'$ comme les éléments du système associé à la « disposition produit » $E \times E'$, avec
\[
  (x, x') \mathrel{|\circ|} (y, y') \overset{\text{déf}}{\iff} x \mathrel{|\circ|} x' \ \text{ou}\ y \mathrel{|\circ|} y' .
\]
\note{la relation est transcrite telle qu'elle est écrite ; on attendrait plutôt $x \mathrel{|\circ|} y$ ou $x' \mathrel{|\circ|} y'$.}
On \uncertain{vérifie} en effet que les \struck{\uncertain{applications}} $R \in \Sigma_{E \times E'}$ sont dans l'ens.\ noté $\operatorname{Corr}^{*}(\underline{E}, \underline{E}')$ (p.~94) i.e.\ \struck{$\forall\,$} $\forall\, x \in E$, $R(x) \in \Sigma_{E'}$, et $\forall\, x' \in E'$, ${}^{t}R(x') \in \Sigma_E$. Mais il ne semble pas que l'inverse soit vrai, \uncertain{sauf} dans le cas où $\underline{E}$ ou $\underline{E}'$ est \uncertain{une} disposition triviale.

Mais si ce n'était pas vrai en général, la question demeure si l'ensemble des \add{ensembles des} \uncertain{dispos.}\ corr.\ de $\underline{E}$ \uncertain{à} $\underline{E}'$, \ldots\ \struck{correspondances} des opérations $\Sigma \to \Sigma'$, forment un système. La relation d'ordre \uncertain{induite} par le cas d'ensembles (disp.\ triviales) \uncertain{\ill{}} \uncertain{l'inclusion}
\note{la page, et la batch, s'arrêtent sur cette phrase ; la suite des « questions ouvertes » annoncées p.~99 se poursuit au-delà.}

\end{document}
